% 核心观点：将一般任务完成函数实例化为同时满足入区与布料复原的视觉可审核事件。
% 叙述逻辑：参考投影与谓词 -> 同时完成条件 -> 可见性和协议边界。
\subsubsection{Bimanual Exposure and Execution}
评价首先在参考视角定义完整目标区域足迹 $R^\star$，
并令 $\Omega$ 为该参考视角的整幅图像域。
令 $O_t^\star$ 为目标完整投影，$C_t$ 为遮挡布料投影。
当目标完整可评估且 $|O_t^\star|>0$ 时，定义
\begin{align}
 p_{\rm in}(t)&=\ind\{O_t^\star\subseteq R^\star\},\\
 a_C(t)&=|C_t\cap\Omega|/|\Omega|,\\
 p_{\rm restore}(t)&=\ind\{a_C(t)\ge\theta_C\}.
\end{align}
$R^\star$ 不是当前可见region类别mask：
互斥分割使object与可见region不重叠，直接求交会误判。
固定目标区域可由无遮挡参考帧独立标注；相机或区域移动时须重新估计。
仅有局部可见的目标mask也不能替代 $O_t^\star$ 来证明整个物体已进入区域。

当前布料复原\emph{不会重新遮挡目标物体}，复原指标是布料在画面中的面积占比，
不是目标被遮挡比例，也不是默认裁剪桌面ROI后的占比。
按这一要求，任务在规定时限内满足下式即可完成：
\begin{equation}
 \Phi(s_{0:T})=
 \ind\{\exists\,t\le T:
 p_{\rm in}(t)=1\ \land\ p_{\rm restore}(t)=1\}.
 \label{eq:completiontask}
\end{equation}
两项必须在同一完成时刻成立：不能将不同时间分别满足、却从未共同成立的记录组合为成功。
式\eqref{eq:completiontask}采用时限内曾共同满足的标准，不额外要求保持至评价终点；
若最终协议要求终态保持或首次完成后停止，须在采集前明确并相应统一评价实现。
若协议另要求连续 $d_C$ 帧确认，则在同一窗口内检验这两个谓词，
并要求窗口终点不超过 $T$；未指定持续时间时不默认添加这一要求。
\draftnote{待补：统一时限 $T$、布料画面覆盖阈值 $\theta_C$ 与持续确认窗口 $d_C$，不沿用未经核实的历史阈值。}

若最终物体与目标区域均可清晰评价，式\eqref{eq:completiontask}的两个谓词
可以由同一参考帧判定，不必依赖假设的最终目标遮挡。
执行中被机械臂临时遮挡仍可能导致评价证据不足，但不能把这种局部情况
混同于任务要求布料最后重新遮住物体。
暴露、分离、推动和复原是有助于解释行为的过程事件；
并非所有任务都必须采用同一组过程标签，更不要求在问题定义中指定阶段头。
